\documentclass[12pt]{article}
\usepackage[utf8]{inputenc}
\usepackage[T1]{fontenc}
\usepackage[margin=1in]{geometry}
\usepackage{setspace}
\usepackage{xcolor}
\usepackage[colorlinks=true, urlcolor=blue!60!black, linkcolor=blue!60!black]{hyperref}
\setlength{\parindent}{0em}
\setlength{\parskip}{0.45em}
\renewcommand{\baselinestretch}{1.0}
\begin{document}
\begin{center}
{\Large \textit{Luther v. Borden}, 48 U.S. (7 How.) 1 (1849)\par}
\vspace{0.3em}
{\normalsize Case Brief --- One-Page Edition\par}
\end{center}

\textbf{Memory Jogger.} The Dorr Rebellion: Guarantee Clause disputes are political questions (\href{https://supreme.justia.com/cases/federal/us/48/1/}{48 U.S.\ 1}).

\textbf{Facts.} Rhode Island's 1663 charter limited voting to freeholders. Reformers adopted an extralegal ``People's Constitution'' and elected Thomas Dorr governor in 1842. The charter government declared martial law, and militiaman Luther Borden broke into Martin Luther's home to arrest him for aiding Dorr.

\textbf{Procedural History.} In Luther's federal trespass suit, Justice Story excluded his evidence of ratification; the jury found for Borden.

\textbf{Votes.} Affirmed, 5--1 (Taney for the Court; Woodbury dissenting in part; three Justices not participating).

\textbf{Holding.} Federal courts cannot decide which rival regime is a state's lawful government or whether it is republican; that is for Congress and the President, whose recognition binds the courts. A state may declare martial law to put down an insurrection.

\textbf{Reasoning.} Voiding the charter government would nullify a year of laws; Rhode Island's courts had upheld it; no manageable standards existed for counting voters; and the Guarantee Clause commits the question to Congress and to the President under the 1795 Act. Woodbury agreed on the political question but rejected martial law over civilians while courts were open.

\textbf{Examples.} \emph{Same-Side:} a Guarantee Clause challenge to ballot initiatives. \emph{Opposite-Side:} Equal Protection malapportionment claims (\textit{Baker v.\ Carr}). \emph{Fence-Sitter:} a Guarantee Clause attack on commandeering.

\textbf{Critique.} Chemerinsky says \textit{Luther} erased the Guarantee Clause; Tara Leigh Grove reads early cases as deference to the political branches.

\textbf{Key Quotation.} ``Under this article of the Constitution, it rests with Congress to decide what government is the established one in a State. For as the United States guarantee to each State a republican government, Congress must necessarily decide what government is established in the State before it can determine whether it is republican or not.'' --- Chief Justice Taney, 48 U.S.\ (7 How.)\ at 42.

\end{document}
